\hypertarget{appendix-bd39c3bf5d64}{}\label{referencestart:appendix-bd39c3bf5d64}
\section{Matrices as linear maps and basis images}\label{appendix-bd39c3bf5d64-matrices-as-linear-maps-and-basis-images}

\hypertarget{appendix-bd39c3bf5d64-matrices-as-linear-maps-and-basis-images}{}

\textbf{Curriculum depth:} Introduced · \textbf{First seen:} Phase7 · \textbf{Latest development:} Phase7

\begin{allmightroledefinition}

\subsection{Definition and mechanism}\label{appendix-bd39c3bf5d64-definition-and-mechanism}

\hypertarget{appendix-bd39c3bf5d64-definition-and-mechanism}{}

A matrix is a rectangular scalar array. After domain, codomain and ordered bases are fixed, its columns are the images of input basis vectors. An m by n matrix maps n input components to m output components; a rectangular matrix does not have an ordinary two-sided inverse.

\end{allmightroledefinition}

\textbf{Reconstruction cue:} Columns (2,1) and (1,3) map components (3,2) to (8,9).

\subsection{Prerequisites}\label{appendix-bd39c3bf5d64-prerequisites}

\hypertarget{appendix-bd39c3bf5d64-prerequisites}{}

\begin{itemize}
\tightlist
\item
  \href{../07-Vision-XR-and-Playable-Worlds-P01/main.pdf\#nameddest=appendix-ce0341f7ea0d}{Basis coordinates and dimension (Vision, XR, and Playable Worlds, Part 1, Reference notes: Basis coordinates and dimension)}
\end{itemize}

\subsection{Used by}\label{appendix-bd39c3bf5d64-used-by}

\hypertarget{appendix-bd39c3bf5d64-used-by}{}

\begin{itemize}
\tightlist
\item
  \href{../07-Vision-XR-and-Playable-Worlds-P01/main.pdf\#nameddest=appendix-ab31ef447dbf}{Matrix multiplication and composition order (Vision, XR, and Playable Worlds, Part 1, Reference notes: Matrix multiplication and composition order)}
\item
  \hyperref[appendix-da478be01d02]{Determinants orientation and volume}
\item
  \hyperref[appendix-0bf0ca73f735]{Linear affine and projective transformations}
\item
  \hyperref[appendix-ffdc1d4e5350]{Matrix storage indexing and strides}
\item
  \hyperref[appendix-26d8da71f6c4]{Quaternion matrix and axis-angle conversion}
\item
  \href{../07-Vision-XR-and-Playable-Worlds-P01/main.pdf\#nameddest=appendix-8572a4a060c1}{Linear systems and fundamental subspaces (Vision, XR, and Playable Worlds, Part 1, Reference notes: Linear systems and fundamental subspaces)} --- Extends this foundation into linear-systems and decomposition analysis.
\item
  \href{../07-Vision-XR-and-Playable-Worlds-P01/main.pdf\#nameddest=appendix-0b3be18411f9}{Eigenvalues and eigenspaces (Vision, XR, and Playable Worlds, Part 1, Reference notes: Eigenvalues and eigenspaces)} --- Extends this foundation into linear-systems and decomposition analysis. \#\# Recall and next depth
\end{itemize}

Reconstruct the example, state its conventions and find an input that violates its preconditions. Deeper algorithms and measured production performance remain separate from the present model.

\subsection{Source and navigation}\label{appendix-bd39c3bf5d64-source-and-navigation}

\hypertarget{appendix-bd39c3bf5d64-source-and-navigation}{}

\hyperref[01-foundations-i-ch018-keep-the-promises-visible-while-moving-between-representations]{Continuous explanation} · \hyperref[appendix-fd090d982a15-phase-7---parts-601-700]{Coverage and source gap} · \hyperref[appendix-a103e7846d21-phase-7---source-analysis-and-verification]{Sources and verification}
